\documentclass[11pt,oneside]{book}

%====================== PACKAGES ======================
\usepackage[utf8]{inputenc}
\usepackage[T1]{fontenc}
\usepackage{mathptmx}              % Times text + math (textbook look)
\usepackage[margin=1in]{geometry}
\usepackage{amsmath,amssymb,amsthm,mathtools}
\usepackage{bm}
\usepackage{enumitem}
\usepackage{xcolor}
\usepackage[most]{tcolorbox}
\usepackage{booktabs}
\usepackage{array}
\usepackage{graphicx}
\usepackage{fancyhdr}
\usepackage[bookmarks,bookmarksnumbered]{hyperref}
\hypersetup{colorlinks=true,linkcolor=blue!55!black,urlcolor=teal,citecolor=blue}

%====================== COLORS ======================
\definecolor{navy}{HTML}{1F3A5F}
\definecolor{steel}{HTML}{2E5A88}
\definecolor{intublue}{HTML}{EAF2FB}
\definecolor{inturule}{HTML}{3E7CB1}
\definecolor{warmbg}{HTML}{FFF6E9}
\definecolor{warmrule}{HTML}{D08B2C}
\definecolor{keybg}{HTML}{EAF6EE}
\definecolor{keyrule}{HTML}{2E8B57}
\definecolor{redbg}{HTML}{FBECEC}
\definecolor{redrule}{HTML}{B23B3B}

%====================== THEOREM ENVIRONMENTS ======================
\theoremstyle{plain}
\newtheorem{theorem}{Theorem}[chapter]
\newtheorem{lemma}[theorem]{Lemma}
\newtheorem{corollary}[theorem]{Corollary}
\newtheorem{proposition}[theorem]{Proposition}
\theoremstyle{definition}
\newtheorem{definition}[theorem]{Definition}
\newtheorem{example}[theorem]{Example}
\newtheorem{exercise}{Exercise}[chapter]
\theoremstyle{remark}
\newtheorem{remark}[theorem]{Remark}

%====================== CALLOUT BOXES ======================
\newtcolorbox{geometry_box}{enhanced,breakable,colback=intublue,colframe=inturule,
  boxrule=0.4pt,left=8pt,right=8pt,top=6pt,bottom=6pt,arc=2pt,
  borderline west={3pt}{0pt}{inturule},fonttitle=\bfseries,title={Geometric interpretation}}
\newtcolorbox{ml_box}{enhanced,breakable,colback=keybg,colframe=keyrule,
  boxrule=0.4pt,left=8pt,right=8pt,top=6pt,bottom=6pt,arc=2pt,
  borderline west={3pt}{0pt}{keyrule},fonttitle=\bfseries,title={Machine-learning interpretation}}
\newtcolorbox{miscon_box}{enhanced,breakable,colback=redbg,colframe=redrule,
  boxrule=0.4pt,left=8pt,right=8pt,top=6pt,bottom=6pt,arc=2pt,
  borderline west={3pt}{0pt}{redrule},fonttitle=\bfseries,title={Common misconception}}
\newtcolorbox{intuition_box}{enhanced,breakable,colback=warmbg,colframe=warmrule,
  boxrule=0.4pt,left=8pt,right=8pt,top=6pt,bottom=6pt,arc=2pt,
  borderline west={3pt}{0pt}{warmrule},fonttitle=\bfseries,title={Motivation \& intuition}}

%====================== MACROS ======================
\newcommand{\R}{\mathbb{R}}
\newcommand{\C}{\mathbb{C}}
\newcommand{\tr}{\operatorname{tr}}
\newcommand{\rank}{\operatorname{rank}}
\newcommand{\col}{\operatorname{col}}
\newcommand{\row}{\operatorname{row}}
\newcommand{\nul}{\operatorname{null}}
\newcommand{\spn}{\operatorname{span}}
\newcommand{\dimm}{\operatorname{dim}}
\newcommand{\norm}[1]{\left\lVert #1 \right\rVert}
\newcommand{\fn}[1]{\norm{#1}_F}
\newcommand{\sn}[1]{\norm{#1}_2}
\newcommand{\T}{^{\mathsf T}}
\newcommand{\inv}{^{-1}}
\newcommand{\inner}[2]{\langle #1,\, #2\rangle}
\newcommand{\bu}{\mathbf{u}}
\newcommand{\bv}{\mathbf{v}}
\newcommand{\bw}{\mathbf{w}}
\newcommand{\bx}{\mathbf{x}}
\newcommand{\by}{\mathbf{y}}
\newcommand{\bz}{\mathbf{z}}
\newcommand{\bq}{\mathbf{q}}
\newcommand{\bp}{\mathbf{p}}
\newcommand{\ba}{\mathbf{a}}
\newcommand{\bb}{\mathbf{b}}
\newcommand{\be}{\mathbf{e}}
\newcommand{\bzero}{\mathbf{0}}
\newcommand{\bmu}{\boldsymbol{\mu}}
\DeclareMathOperator*{\argmin}{arg\,min}
\DeclareMathOperator*{\argmax}{arg\,max}

\renewcommand{\qedsymbol}{$\blacksquare$}

%====================== HEADERS ======================
\pagestyle{fancy}
\fancyhf{}
\fancyhead[L]{\small\nouppercase{\leftmark}}
\fancyhead[R]{\small\thepage}
\renewcommand{\headrulewidth}{0.4pt}

%====================== TITLE ======================
\title{\vspace{-2cm}
  {\color{navy}\Huge\bfseries The Singular Value Decomposition}\\[10pt]
  {\color{steel}\LARGE A Rigorous Development from First Principles}\\[8pt]
  {\color{steel}\large culminating in the Eckart--Young--Mirsky Theorem}}
\author{\Large A self-contained graduate treatment\\[4pt]
  \normalsize every theorem motivated, every step proved}
\date{\today}

\begin{document}

\frontmatter
\maketitle

\chapter*{Preface}
\addcontentsline{toc}{chapter}{Preface}
\markboth{Preface}{Preface}

This book develops the Singular Value Decomposition (SVD) and its central optimality
property --- the \emph{Eckart--Young--Mirsky best rank-$k$ approximation theorem} ---
in the rigorous, build-everything-from-the-ground-up style of a graduate linear-algebra
text. The guiding principle is simple and uncompromising:

\begin{center}
\emph{Nothing is asserted that has not first been proved.}
\end{center}

We deliberately avoid the phrases ``it is easy to see,'' ``clearly,'' and ``it can be
shown.'' When a theorem depends on another result, that result is stated and proved
\emph{first}. The reader is assumed to know only the mechanics of matrix
multiplication, transpose, inverse, and determinants. Every other object --- span,
basis, rank, inner product, orthogonal projection, symmetric matrices, the Spectral
Theorem, positive semidefiniteness, matrix norms --- is introduced from scratch and
proved before it is used.

Each chapter closes with a \textbf{summary} and a set of \textbf{exercises}. Throughout,
shaded boxes separate the formal mathematics from \emph{motivation}, \emph{geometric
interpretation}, \emph{machine-learning interpretation}, and \emph{common
misconceptions}, so that intuition supports the proofs without ever substituting for
them.

The logical dependency of the chapters is strictly linear:
\begin{gather*}
\text{rank}\ \to\ \text{orthogonality}\ \to\ \text{symmetric matrices}\ \to\
\text{Spectral Theorem}\\[2pt]
\to\ \text{PSD}\ \to\ \text{SVD}\ \to\ \text{norms}\ \to\
\text{Eckart--Young--Mirsky}.
\end{gather*}
Read it in order; each chapter is a prerequisite for the next.

\tableofcontents

\mainmatter

%##################################################################
\chapter{Motivation: Why the SVD Exists}
%##################################################################

\begin{intuition_box}
Eigenvalue decomposition is the most powerful tool in elementary linear algebra: it
turns the action of a matrix into independent scalar stretches along special
directions. But it has two fatal restrictions --- it needs the matrix to be
\emph{square}, and even then it may fail to exist (the matrix may not be
diagonalizable). Data matrices are almost never square, and we cannot afford a tool
that sometimes does not exist. The SVD is the repair: a decomposition that exists for
\emph{every} matrix, square or rectangular, and that recovers all the computational
power of the eigenbasis.
\end{intuition_box}

\section{The eigenvalue decomposition and its limits}

For a square matrix $A\in\R^{n\times n}$ the eigenvalue equation
$A\bv=\lambda\bv$ identifies directions $\bv$ along which $A$ acts as pure scaling. When
$A$ possesses $n$ linearly independent eigenvectors, stacking them as the columns of
$V$ gives the \emph{eigenvalue decomposition} (EVD)
\[
A=V\Lambda V\inv,\qquad \Lambda=\operatorname{diag}(\lambda_1,\dots,\lambda_n),
\]
which converts repeated multiplication by $A$ into repeated scalar multiplication. This
is the engine behind iterative methods, differential equations, and Markov chains.

The decomposition demands two things, and \emph{both} can fail:

\begin{enumerate}[label=(\roman*),leftmargin=2.2em]
\item \textbf{Squareness.} The equation $A\bv=\lambda\bv$ requires $A\bv$ and $\bv$ to
lie in the same space. If $A\in\R^{m\times n}$ with $m\ne n$, then $\bv\in\R^n$ but
$A\bv\in\R^m$, and the equation is not even well posed. Rectangular matrices simply
\emph{cannot} have eigenvectors.
\item \textbf{Diagonalizability.} Even square matrices may lack a full set of
independent eigenvectors. The matrix $\begin{psmallmatrix}1&1\\0&1\end{psmallmatrix}$
has the single eigenvalue $1$ but only a one-dimensional eigenspace, so no
$V\inv$ exists.
\end{enumerate}

\section{What the SVD promises}

The SVD circumvents both problems by abandoning the demand for a \emph{single} basis in
which $A$ is diagonal. Instead it allows \emph{two} orthonormal bases --- one in the
input space $\R^n$, one in the output space $\R^m$ --- linked so that $A$ maps one to
the other with pure scaling in between. We will prove (Chapter~\ref{ch:svd}) that
\emph{every} real matrix $A\in\R^{m\times n}$ factors as
\[
A=U\Sigma V\T,
\]
with $U\in\R^{m\times m}$ and $V\in\R^{n\times n}$ orthogonal and $\Sigma$ a
nonnegative diagonal matrix. No squareness, no diagonalizability hypothesis --- the
factorization always exists.

\section{Why anyone outside pure mathematics should care}

The SVD is arguably the single most useful matrix factorization in applied science.
A non-exhaustive list of where its optimality (the subject of Chapter~\ref{ch:ey})
is doing the work:

\begin{description}[style=nextline,leftmargin=1.6em]
\item[Principal Component Analysis (PCA).] The directions of maximal variance in a data
cloud are exactly the right singular vectors of the centred data matrix
(Chapter~\ref{ch:intuition}). Dimensionality reduction \emph{is} truncated SVD.
\item[Image compression.] A grayscale image is a matrix of pixel intensities. Because
its singular values decay rapidly, a few dozen rank-one layers reproduce the picture
while storing a small fraction of the numbers.
\item[Latent Semantic Analysis (LSA).] A term--document matrix, truncated by SVD,
exposes latent ``topics'' as singular directions, mapping synonyms to nearby vectors.
\item[Recommendation systems.] A user--item rating matrix is modelled as approximately
low rank; the missing entries are predicted by its best low-rank approximation.
\item[LoRA (Low-Rank Adaptation).] Fine-tuning updates to enormous neural networks are
empirically near low rank, so they are stored in factored form $\Delta W=BA$ with $B,A$
thin --- a direct application of the low-rank idea, slashing trainable parameters by
orders of magnitude.
\item[Numerical stability.] The SVD gives the most reliable computation of rank,
pseudoinverse, and least-squares solutions, precisely because it never forms the
condition-number-squaring product $A\T A$.
\end{description}

Each of these rests on the theorem at the heart of this book: among all matrices of
rank at most $k$, the truncated SVD is provably the closest to $A$. Everything we build
in Chapters~\ref{ch:rank}--\ref{ch:norms} is in service of proving that statement
cleanly in Chapter~\ref{ch:ey}.

\section*{Summary}
\addcontentsline{toc}{section}{Summary}
EVD is powerful but restricted to square, diagonalizable matrices. The SVD removes both
restrictions by using paired orthonormal bases in the input and output spaces. It exists
universally and underlies PCA, compression, LSA, recommender systems, and LoRA. The
remainder of the book proves the SVD exists and that its truncation is optimal.

\begin{exercise}
Exhibit a $2\times2$ real matrix that has no real eigenvalues, and verify directly that
it nonetheless has a real SVD (you may use a computer for the arithmetic; we prove
existence in general later).
\end{exercise}
\begin{exercise}
Explain in one paragraph why the equation $A\bv=\lambda\bv$ is dimensionally
inconsistent when $A\in\R^{3\times 5}$.
\end{exercise}
\begin{exercise}
The matrix $\begin{psmallmatrix}2&1\\0&2\end{psmallmatrix}$ is square. Show it has only a
one-dimensional eigenspace, hence no EVD of the form $V\Lambda V\inv$ with diagonal
$\Lambda$ exists.
\end{exercise}

%##################################################################
\chapter{Linear-Algebra Prerequisites: Rank}\label{ch:rank}
%##################################################################

We develop only what the SVD requires, but we develop it completely. Throughout,
vectors are columns and $\R^n$ carries the usual operations.

\section{Linear independence, span, basis, dimension}

\begin{definition}[Linear combination, span]
A \emph{linear combination} of vectors $\bv_1,\dots,\bv_k\in\R^n$ is any vector
$\sum_{i=1}^k c_i\bv_i$ with $c_i\in\R$. Their \emph{span} is the set of all such
combinations,
$\spn\{\bv_1,\dots,\bv_k\}=\{\sum_i c_i\bv_i: c_i\in\R\}$,
which is a subspace of $\R^n$.
\end{definition}

\begin{definition}[Linear independence]
The vectors $\bv_1,\dots,\bv_k$ are \emph{linearly independent} if the only scalars
solving $\sum_{i=1}^k c_i\bv_i=\bzero$ are $c_1=\cdots=c_k=0$. Otherwise they are
\emph{linearly dependent}.
\end{definition}

\begin{definition}[Basis, dimension]
A \emph{basis} of a subspace $W\subseteq\R^n$ is a linearly independent set that spans
$W$. The number of vectors in a basis is the \emph{dimension} $\dimm W$; we take as
known from elementary linear algebra that every basis of a given subspace has the same
size, so $\dimm W$ is well defined.
\end{definition}

\begin{lemma}[Coordinates are unique in a basis]\label{lem:uniquecoord}
If $\bv_1,\dots,\bv_k$ is a basis of $W$ then every $\bw\in W$ has a \emph{unique}
representation $\bw=\sum_i c_i\bv_i$.
\end{lemma}
\begin{proof}
Existence holds because the basis spans $W$. For uniqueness, suppose
$\sum_i c_i\bv_i=\sum_i d_i\bv_i$. Subtracting gives $\sum_i(c_i-d_i)\bv_i=\bzero$, and
linear independence forces $c_i-d_i=0$ for all $i$, i.e.\ $c_i=d_i$.
\end{proof}

\section{The four fundamental subspaces and rank}

\begin{definition}[Column space, row space, null space]
For $A\in\R^{m\times n}$:
\begin{itemize}[leftmargin=1.5em]
\item the \emph{column space} $\col(A)=\{A\bx:\bx\in\R^n\}\subseteq\R^m$ is the span of
the columns of $A$;
\item the \emph{row space} $\row(A)=\col(A\T)\subseteq\R^n$ is the span of the rows;
\item the \emph{null space} $\nul(A)=\{\bx\in\R^n:A\bx=\bzero\}\subseteq\R^n$.
\end{itemize}
\end{definition}

\begin{definition}[Rank: three candidate definitions]\label{def:rank}
For $A\in\R^{m\times n}$ define
\[
r_{\mathrm{col}}=\dimm\col(A),\qquad
r_{\mathrm{row}}=\dimm\row(A),\qquad
r_{\mathrm{fac}}=\min\{k: A=YZ\T,\ Y\in\R^{m\times k},\ Z\in\R^{n\times k}\}.
\]
\end{definition}

These three numbers coincide. We prove it, because the equality (``row rank $=$ column
rank'') is used repeatedly later and should not be taken on faith.

\begin{theorem}[Equivalence of the rank definitions]\label{thm:rankequiv}
For every $A\in\R^{m\times n}$, $r_{\mathrm{col}}=r_{\mathrm{row}}=r_{\mathrm{fac}}$.
The common value is the \emph{rank} $\rank(A)$.
\end{theorem}
\begin{proof}
\textbf{Step 1: $r_{\mathrm{col}}=r_{\mathrm{fac}}$.}
Let $r=r_{\mathrm{col}}$ and pick a basis $\bb_1,\dots,\bb_r$ of $\col(A)$; assemble it
as the columns of $Y=[\bb_1\ \cdots\ \bb_r]\in\R^{m\times r}$. Each column $\ba_j$ of
$A$ lies in $\col(A)$, so by Lemma~\ref{lem:uniquecoord} it has coordinates
$\ba_j=\sum_{i=1}^r z_{ji}\bb_i=Y\bz_j$ where $\bz_j=(z_{j1},\dots,z_{jr})\T$. Stacking
the $\bz_j$ as rows of a matrix $Z\in\R^{n\times r}$ gives $A=YZ\T$. Hence a
factorization of inner dimension $r$ exists, so $r_{\mathrm{fac}}\le r$. Conversely, if
$A=YZ\T$ with $Y\in\R^{m\times k}$ then every column of $A$ is
$A\be_j=Y(Z\T\be_j)\in\col(Y)$, so $\col(A)\subseteq\col(Y)$ and
$r_{\mathrm{col}}\le\dimm\col(Y)\le k$. Taking $k=r_{\mathrm{fac}}$ yields
$r_{\mathrm{col}}\le r_{\mathrm{fac}}$. Combining, $r_{\mathrm{col}}=r_{\mathrm{fac}}$.

\textbf{Step 2: $r_{\mathrm{row}}=r_{\mathrm{fac}}$.}
Apply Step 1 to $A\T$: its column space is $\row(A)$, and $A\T=Z Y\T$ is a
factorization of inner dimension $r_{\mathrm{fac}}$ (the same $Y,Z$ as above,
transposed). Therefore $r_{\mathrm{row}}=\dimm\col(A\T)$ equals the minimal inner
dimension of a factorization of $A\T$, which is $r_{\mathrm{fac}}$ because transposing a
factorization of $A$ gives one of $A\T$ of the same inner dimension and vice versa.

Hence all three numbers agree.
\end{proof}

\begin{remark}
Theorem~\ref{thm:rankequiv} is the statement, familiar from a first course but rarely
proved there, that \emph{the maximal number of independent columns equals the maximal
number of independent rows}. We will use the factorization form
$A=YZ\T$ (Figure~``long-skinny times short-long'') constantly.
\end{remark}

\section{Rank-$0$, rank-$1$, rank-$2$, rank-$k$}

\begin{definition}[Rank-$k$ matrix]
$A$ has rank $k$ if $\rank(A)=k$. Equivalently (Theorem~\ref{thm:rankequiv}) $A$ is a
sum of $k$ rank-one matrices but not of $k-1$.
\end{definition}

The only rank-$0$ matrix is the zero matrix. Rank-one matrices have an exceptionally
clean description, which is the workhorse of the entire theory.

\begin{theorem}[Characterization of rank-one matrices]\label{thm:rank1}
A matrix $A\in\R^{m\times n}$ has rank $1$ if and only if $A=\bu\bv\T$ for some nonzero
vectors $\bu\in\R^m$, $\bv\in\R^n$.
\end{theorem}
\begin{proof}
($\Leftarrow$) Suppose $A=\bu\bv\T$ with $\bu,\bv\ne\bzero$. Column $j$ of $A$ is
$A\be_j=\bu(\bv\T\be_j)=v_j\,\bu$, a scalar multiple of the fixed nonzero vector $\bu$.
Thus every column lies in $\spn\{\bu\}$, so $\col(A)\subseteq\spn\{\bu\}$, giving
$\dimm\col(A)\le 1$. Since some $v_j\ne 0$, the column $v_j\bu\ne\bzero$, so $A\ne 0$ and
$\dimm\col(A)\ge 1$. Hence $\rank(A)=1$.

($\Rightarrow$) Suppose $\rank(A)=1$. Then $\col(A)$ is one-dimensional; let $\bu$ be a
nonzero spanning vector, so $\col(A)=\spn\{\bu\}$. Each column $\ba_j\in\col(A)$ is a
multiple $\ba_j=v_j\bu$ for a unique scalar $v_j$ (Lemma~\ref{lem:uniquecoord}). Setting
$\bv=(v_1,\dots,v_n)\T$ gives, column by column, $A=[\,v_1\bu\ \cdots\ v_n\bu\,]
=\bu\bv\T$. As $A\ne 0$, $\bv\ne\bzero$.
\end{proof}

\begin{example}
$A=\begin{psmallmatrix}2&4&6\\1&2&3\end{psmallmatrix}$ has every row a multiple of
$(2,4,6)$ and every column a multiple of $(2,1)\T$. Indeed
$A=\begin{psmallmatrix}2\\1\end{psmallmatrix}\begin{psmallmatrix}1&2&3\end{psmallmatrix}
=\bu\bv\T$, confirming $\rank(A)=1$.
\end{example}

\begin{remark}[Rank-two and rank-$k$ as superpositions]
A rank-two matrix is a sum $\bu_1\bv_1\T+\bu_2\bv_2\T$ that is not itself rank $0$ or
$1$. In general, Theorem~\ref{thm:rankequiv} says rank-$k$ matrices are exactly the sums
of $k$ rank-one matrices that cannot be written with fewer. The SVD
(Chapter~\ref{ch:svd}) will hand us a \emph{canonical} such decomposition in which the
rank-one pieces are mutually orthogonal and ordered by importance.
\end{remark}

\begin{ml_box}
Low rank is the mathematical face of \emph{structure} or \emph{redundancy} in data. A
rank-one matrix $\bu\bv\T$ is the cleanest possible structure: one ``pattern'' $\bu$
modulated across columns by the weights in $\bv$. When a data matrix is approximately a
sum of a few rank-one pieces, it is approximately low rank, and the whole apparatus of
low-rank approximation (compression, denoising, completion) becomes available.
\end{ml_box}

\section*{Summary}
\addcontentsline{toc}{section}{Summary}
Span, independence, basis, and dimension give meaning to the four fundamental subspaces.
Rank has three equivalent definitions --- column rank, row rank, and minimal
factorization dimension --- all proved equal in Theorem~\ref{thm:rankequiv}. Rank-one
matrices are exactly the outer products $\bu\bv\T$ (Theorem~\ref{thm:rank1}), and
rank-$k$ matrices are minimal sums of $k$ such products.

\begin{exercise}
Prove that $\rank(AB)\le\min\{\rank(A),\rank(B)\}$ using the factorization
characterization $r_{\mathrm{fac}}$.
\end{exercise}
\begin{exercise}
Show that for $A=\bu\bv\T$ (nonzero) the null space has dimension $n-1$, and identify it
explicitly in terms of $\bv$.
\end{exercise}
\begin{exercise}
Give an example of $2\times2$ matrices with $\rank(A)=\rank(B)=1$ but $\rank(A+B)=2$,
and another pair with $\rank(A+B)=1$. Conclude that rank is neither additive nor
subadditive in an exact sense (only the inequality $\rank(A+B)\le\rank A+\rank B$ holds).
\end{exercise}
\begin{exercise}
Prove the inequality $\rank(A+B)\le\rank(A)+\rank(B)$ from Theorem~\ref{thm:rankequiv}.
\end{exercise}

%##################################################################
\chapter{Orthogonality}\label{ch:orth}
%##################################################################

\begin{intuition_box}
Orthogonality is what makes the SVD \emph{clean}. Without it, decomposing a matrix into
ingredients would give pieces that overlap and interfere; with it, the pieces are
mutually perpendicular, energies simply add (Pythagoras), and projecting onto a subspace
is the optimal way to approximate. Every optimality statement in this book ultimately
rests on the geometry of right angles.
\end{intuition_box}

\section{Inner product, norm, orthogonality}

\begin{definition}[Inner product and norm]
For $\bx,\by\in\R^n$ the \emph{inner (dot) product} is $\inner{\bx}{\by}=\bx\T\by=
\sum_{i=1}^n x_iy_i$, and the \emph{Euclidean norm} is
$\norm{\bx}=\sqrt{\inner{\bx}{\bx}}=\bigl(\sum_i x_i^2\bigr)^{1/2}$. Vectors $\bx,\by$
are \emph{orthogonal}, written $\bx\perp\by$, if $\inner{\bx}{\by}=0$.
\end{definition}

\begin{proposition}[Basic properties of the inner product]\label{prop:innerprops}
For all $\bx,\by,\bz\in\R^n$ and $c\in\R$:
\textup{(i)} symmetry $\inner{\bx}{\by}=\inner{\by}{\bx}$;
\textup{(ii)} linearity $\inner{c\bx+\bz}{\by}=c\inner{\bx}{\by}+\inner{\bz}{\by}$;
\textup{(iii)} positivity $\inner{\bx}{\bx}\ge0$ with equality iff $\bx=\bzero$;
\textup{(iv)} the adjoint identity $\inner{A\bx}{\by}=\inner{\bx}{A\T\by}$ for any
$A\in\R^{m\times n}$.
\end{proposition}
\begin{proof}
(i)--(iii) are immediate from $\inner{\bx}{\by}=\sum_i x_iy_i$. For (iv),
$\inner{A\bx}{\by}=(A\bx)\T\by=\bx\T A\T\by=\inner{\bx}{A\T\by}$, using
$(A\bx)\T=\bx\T A\T$.
\end{proof}

\begin{theorem}[Cauchy--Schwarz inequality]\label{thm:cs}
For all $\bx,\by\in\R^n$, $\ |\inner{\bx}{\by}|\le\norm{\bx}\,\norm{\by}$, with equality
iff $\bx,\by$ are linearly dependent.
\end{theorem}
\begin{proof}
If $\by=\bzero$ both sides are $0$. Otherwise, for every $t\in\R$ positivity gives
\[
0\le\norm{\bx-t\by}^2=\inner{\bx-t\by}{\bx-t\by}
=\norm{\bx}^2-2t\inner{\bx}{\by}+t^2\norm{\by}^2 =: q(t).
\]
The quadratic $q(t)$ is nonnegative for all $t$, so its discriminant is $\le0$:
$4\inner{\bx}{\by}^2-4\norm{\bx}^2\norm{\by}^2\le0$, i.e.\
$\inner{\bx}{\by}^2\le\norm{\bx}^2\norm{\by}^2$. Taking square roots gives the
inequality. Equality forces $q(t_0)=0$ for the vertex $t_0$, i.e.\ $\bx=t_0\by$,
linear dependence; the converse is a direct check.
\end{proof}

\begin{theorem}[Pythagorean theorem]\label{thm:pyth}
If $\bx\perp\by$ then $\norm{\bx+\by}^2=\norm{\bx}^2+\norm{\by}^2$.
\end{theorem}
\begin{proof}
$\norm{\bx+\by}^2=\inner{\bx+\by}{\bx+\by}=\norm{\bx}^2+2\inner{\bx}{\by}+\norm{\by}^2$,
and the middle term vanishes because $\inner{\bx}{\by}=0$.
\end{proof}

\section{Orthonormal sets and orthogonal matrices}

\begin{definition}[Orthonormal set]
Vectors $\bq_1,\dots,\bq_k$ are \emph{orthonormal} if $\inner{\bq_i}{\bq_j}=\delta_{ij}$
(unit length, mutually orthogonal).
\end{definition}

\begin{proposition}[Orthonormal $\Rightarrow$ independent]\label{prop:orthoindep}
An orthonormal set is linearly independent.
\end{proposition}
\begin{proof}
Suppose $\sum_i c_i\bq_i=\bzero$. Take the inner product with $\bq_j$:
$0=\inner{\bzero}{\bq_j}=\sum_i c_i\inner{\bq_i}{\bq_j}=\sum_i c_i\delta_{ij}=c_j$. Thus
every $c_j=0$.
\end{proof}

\begin{definition}[Orthogonal matrix]
A square matrix $Q\in\R^{n\times n}$ is \emph{orthogonal} if $Q\T Q=I$. Equivalently its
columns are orthonormal.
\end{definition}

\begin{proposition}[Properties of orthogonal matrices]\label{prop:Qprops}
Let $Q\in\R^{n\times n}$ be orthogonal. Then
\textup{(i)} $Q$ is invertible with $Q\inv=Q\T$, hence also $QQ\T=I$ and the rows are
orthonormal;
\textup{(ii)} $Q$ preserves inner products: $\inner{Q\bx}{Q\by}=\inner{\bx}{\by}$;
\textup{(iii)} $Q$ preserves norms: $\norm{Q\bx}=\norm{\bx}$;
\textup{(iv)} the product of orthogonal matrices is orthogonal.
\end{proposition}
\begin{proof}
(i) $Q\T Q=I$ exhibits $Q\T$ as a left inverse of a square matrix, which is therefore
the two-sided inverse, so $QQ\T=I$ as well; reading $QQ\T=I$ columnwise says the rows of
$Q$ are orthonormal. (ii) $\inner{Q\bx}{Q\by}=(Q\bx)\T(Q\by)=\bx\T Q\T Q\by=\bx\T\by
=\inner{\bx}{\by}$. (iii) Put $\by=\bx$ in (ii) and take square roots. (iv) If
$Q_1\T Q_1=I=Q_2\T Q_2$ then $(Q_1Q_2)\T(Q_1Q_2)=Q_2\T Q_1\T Q_1Q_2=Q_2\T Q_2=I$.
\end{proof}

\begin{geometry_box}
An orthogonal matrix is a \emph{rigid motion} of space: a rotation, a reflection, or a
combination. It moves points around without ever stretching, shrinking, or shearing.
Property (iii) is the precise statement that lengths are untouched, and (ii) that angles
are untouched. This is why orthogonal factors in the SVD carry no ``size'' information ---
all the scaling lives in the diagonal $\Sigma$.
\end{geometry_box}

\section{Orthogonal projection onto a subspace}

\begin{definition}[Orthogonal projector]
Let $W\subseteq\R^m$ be a subspace with orthonormal basis $\bq_1,\dots,\bq_d$, collected
as columns of $Q\in\R^{m\times d}$ (so $Q\T Q=I_d$). The \emph{orthogonal projector}
onto $W$ is $P=QQ\T=\sum_{j=1}^d\bq_j\bq_j\T\in\R^{m\times m}$.
\end{definition}

\begin{proposition}[Defining properties of a projector]\label{prop:projprops}
With $P=QQ\T$ as above: \textup{(i)} $P\T=P$ (symmetric); \textup{(ii)} $P^2=P$
(idempotent); \textup{(iii)} $P\bw=\bw$ for $\bw\in W$ and $P\bx=\bzero$ for
$\bx\perp W$; \textup{(iv)} $P$ does not depend on the choice of orthonormal basis of
$W$.
\end{proposition}
\begin{proof}
(i) $(QQ\T)\T=QQ\T$. (ii) $P^2=QQ\T QQ\T=Q(Q\T Q)Q\T=QI_dQ\T=QQ\T=P$. (iii) For
$\bw=Q\mathbf{c}\in W$, $P\bw=QQ\T Q\mathbf{c}=Q\mathbf{c}=\bw$; for $\bx\perp W$ each
$\bq_j\T\bx=0$, so $Q\T\bx=\bzero$ and $P\bx=\bzero$. (iv) Any $\bx\in\R^m$ splits as
$\bx=\bw+\bx^\perp$ with $\bw\in W$, $\bx^\perp\perp W$ (existence of such a split is the
orthogonal decomposition, shown next in Theorem~\ref{thm:bestapprox}); by (iii)
$P\bx=\bw$, which is determined by $W$ and $\bx$ alone, not by the basis.
\end{proof}

\begin{theorem}[Best approximation / orthogonal decomposition]\label{thm:bestapprox}
Let $W\subseteq\R^m$ be a subspace with orthogonal projector $P$. For every
$\by\in\R^m$:
\begin{enumerate}[label=\textup{(\alph*)},leftmargin=2em]
\item $\by-P\by\perp W$, so $\by=P\by+(\by-P\by)$ is an orthogonal decomposition;
\item for every $\bw\in W$, $\ \norm{\by-\bw}^2=\norm{\by-P\by}^2+\norm{P\by-\bw}^2$;
\item consequently $\displaystyle\min_{\bw\in W}\norm{\by-\bw}=\norm{\by-P\by}$, with the
unique minimizer $\bw=P\by$.
\end{enumerate}
\end{theorem}
\begin{proof}
(a) For each basis vector $\bq_j$ of $W$,
$\inner{\by-P\by}{\bq_j}=\inner{\by}{\bq_j}-\inner{QQ\T\by}{\bq_j}
=\inner{\by}{\bq_j}-(Q\T\by)\T(Q\T\bq_j)$. Now $Q\T\bq_j=\be_j$, so
$(Q\T\by)\T\be_j=(Q\T\by)_j=\bq_j\T\by=\inner{\by}{\bq_j}$, and the difference is $0$.
Hence $\by-P\by$ is orthogonal to every basis vector, thus to all of $W$.
(b) Write $\by-\bw=(\by-P\by)+(P\by-\bw)$. The first summand is $\perp W$ by (a), and
the second lies in $W$; they are orthogonal, so Pythagoras
(Theorem~\ref{thm:pyth}) gives the displayed identity.
(c) In (b) the first term is fixed and the second is $\ge0$, minimized (uniquely) when
$P\by-\bw=\bzero$, i.e.\ $\bw=P\by$.
\end{proof}

\begin{corollary}[Matrix Pythagoras]\label{cor:matpyth}
For any orthogonal projector $P\in\R^{m\times m}$ and any $A\in\R^{m\times n}$,
\[
\fn{A}^2=\fn{PA}^2+\fn{(I-P)A}^2 .
\]
\end{corollary}
\begin{proof}
Apply Theorem~\ref{thm:bestapprox}(b) with $\bw=P\ba_j$ to each column $\ba_j$ of $A$:
$\norm{\ba_j}^2=\norm{P\ba_j}^2+\norm{\ba_j-P\ba_j}^2$. The Frobenius norm squared is the
sum of squared column norms (proved in Chapter~\ref{ch:norms}, but here we only use the
definition $\fn{A}^2=\sum_j\norm{\ba_j}^2$), so summing over $j$ gives the claim, noting
$(I-P)\ba_j=\ba_j-P\ba_j$.
\end{proof}

\begin{miscon_box}
``Projection just means dropping the last coordinates.'' That is only true when $W$ is a
coordinate subspace and the basis is the standard one. In general the projector is
$P=QQ\T$ for an orthonormal basis $Q$ of $W$; it is the \emph{nearest-point map} onto
$W$, and the direction ``dropped'' is the orthogonal complement of $W$, not any fixed
set of axes.
\end{miscon_box}

\section*{Summary}
\addcontentsline{toc}{section}{Summary}
Inner products give length, angle, Cauchy--Schwarz (Theorem~\ref{thm:cs}), and Pythagoras
(Theorem~\ref{thm:pyth}). Orthonormal sets are automatically independent
(Proposition~\ref{prop:orthoindep}); orthogonal matrices preserve all geometry
(Proposition~\ref{prop:Qprops}). The orthogonal projector $P=QQ\T$ realizes the
nearest-point map (Theorem~\ref{thm:bestapprox}), and the matrix Pythagoras identity
(Corollary~\ref{cor:matpyth}) will drive the optimality proof in Chapter~\ref{ch:ey}.

\begin{exercise}
Prove the triangle inequality $\norm{\bx+\by}\le\norm{\bx}+\norm{\by}$ from
Cauchy--Schwarz.
\end{exercise}
\begin{exercise}
Show that a square matrix preserves all norms (i.e.\ $\norm{Q\bx}=\norm{\bx}$ for all
$\bx$) if and only if it is orthogonal. \emph{Hint:} polarization recovers the inner
product from the norm.
\end{exercise}
\begin{exercise}
Let $P$ be symmetric and idempotent ($P\T=P=P^2$). Prove $P$ is the orthogonal projector
onto $\col(P)$, and that $I-P$ is the projector onto $\col(P)^\perp$.
\end{exercise}
\begin{exercise}
Compute the projector onto $W=\spn\{(1,1,0)\T,(0,1,1)\T\}\subseteq\R^3$ by first
orthonormalizing the basis (Gram--Schmidt), then forming $QQ\T$.
\end{exercise}

%##################################################################
\chapter{Symmetric Matrices}\label{ch:sym}
%##################################################################

\begin{intuition_box}
The SVD will be manufactured out of the symmetric matrices $A\T A$ and $AA\T$. Symmetric
matrices are the best-behaved objects in linear algebra: their eigenvalues are real,
their eigenvectors can be chosen orthonormal, and they are diagonalized by an
\emph{orthogonal} change of basis. This chapter proves the two facts we need and the
next chapter assembles them into the Spectral Theorem.
\end{intuition_box}

\begin{definition}[Symmetric matrix]
$A\in\R^{n\times n}$ is \emph{symmetric} if $A\T=A$, i.e.\ $a_{ij}=a_{ji}$ for all
$i,j$.
\end{definition}

\begin{lemma}[Symmetry is self-adjointness]\label{lem:selfadj}
$A$ is symmetric iff $\inner{A\bx}{\by}=\inner{\bx}{A\by}$ for all $\bx,\by\in\R^n$.
\end{lemma}
\begin{proof}
By Proposition~\ref{prop:innerprops}(iv), $\inner{A\bx}{\by}=\inner{\bx}{A\T\by}$ always.
If $A=A\T$ this equals $\inner{\bx}{A\by}$. Conversely, if
$\inner{\bx}{A\T\by}=\inner{\bx}{A\by}$ for all $\bx,\by$ then
$\inner{\bx}{(A\T-A)\by}=0$ for all $\bx$, forcing $(A\T-A)\by=\bzero$ for all $\by$,
hence $A\T=A$.
\end{proof}

\begin{theorem}[Real eigenvalues]\label{thm:realeig}
Every eigenvalue of a real symmetric matrix $A$ is real, and a corresponding eigenvector
may be taken real.
\end{theorem}
\begin{proof}
Allow complex vectors and use the Hermitian inner product
$\inner{\bx}{\by}_{\C}=\bar\bx\T\by$. Suppose $A\bz=\lambda\bz$ with $\bz\ne\bzero$,
$\lambda\in\C$. Then
\[
\bar\bz\T A\bz=\lambda\,\bar\bz\T\bz=\lambda\norm{\bz}^2 .
\]
Taking the conjugate transpose of the scalar $\bar\bz\T A\bz$ and using $A$ real
symmetric ($\bar A=A$, $A\T=A$),
\[
\overline{\bar\bz\T A\bz}=\bz\T A\bar\bz=(\,\overline{A\bz}\,)\T\bz
=\overline{\lambda}\,\bar\bz\T\bz=\overline{\lambda}\norm{\bz}^2 .
\]
A scalar equals its own conjugate, so $\lambda\norm{\bz}^2=\overline\lambda\norm{\bz}^2$;
since $\norm{\bz}^2>0$, $\lambda=\overline\lambda$, i.e.\ $\lambda\in\R$. With $\lambda$
real, $(A-\lambda I)$ is a real singular matrix, so it has a real null vector, providing
a real eigenvector.
\end{proof}

\begin{theorem}[Orthogonality of eigenvectors for distinct eigenvalues]\label{thm:eigorth}
If $A$ is symmetric and $A\bx=\lambda\bx$, $A\by=\mu\by$ with $\lambda\ne\mu$, then
$\bx\perp\by$.
\end{theorem}
\begin{proof}
Using Lemma~\ref{lem:selfadj} two ways,
\[
\lambda\inner{\bx}{\by}=\inner{\lambda\bx}{\by}=\inner{A\bx}{\by}
=\inner{\bx}{A\by}=\inner{\bx}{\mu\by}=\mu\inner{\bx}{\by}.
\]
Hence $(\lambda-\mu)\inner{\bx}{\by}=0$. As $\lambda\ne\mu$, $\inner{\bx}{\by}=0$.
\end{proof}

\begin{remark}
Theorem~\ref{thm:eigorth} handles \emph{distinct} eigenvalues. When an eigenvalue is
repeated, its eigenspace is multi-dimensional and we are free to \emph{choose} an
orthonormal basis within it (Gram--Schmidt). The Spectral Theorem of the next chapter
packages both cases into a single orthonormal eigenbasis of all of $\R^n$.
\end{remark}

\section*{Summary}
\addcontentsline{toc}{section}{Summary}
Symmetry equals self-adjointness (Lemma~\ref{lem:selfadj}). Real symmetric matrices have
real eigenvalues (Theorem~\ref{thm:realeig}) and orthogonal eigenvectors across distinct
eigenvalues (Theorem~\ref{thm:eigorth}). These are exactly the ingredients the Spectral
Theorem needs.

\begin{exercise}
Show that if $A$ is symmetric then $\bx\T A\bx$ is real for all real $\bx$ (trivially),
and that $\bx\T A\bx=\sum_{i}\lambda_i(\bq_i\T\bx)^2$ once the Spectral Theorem is
available.
\end{exercise}
\begin{exercise}
Give a $2\times2$ \emph{non}-symmetric matrix with complex eigenvalues, illustrating that
Theorem~\ref{thm:realeig} genuinely uses symmetry.
\end{exercise}
\begin{exercise}
Prove that for symmetric $A$, eigenvectors from a repeated eigenvalue need not be
orthogonal as given, but can always be \emph{re}chosen orthonormal.
\end{exercise}

%##################################################################
\chapter{The Spectral Theorem}\label{ch:spectral}
%##################################################################

\begin{intuition_box}
The Spectral Theorem is the climax of the symmetric-matrix story and the foundation of
the SVD. It says a symmetric matrix is, after an orthogonal change of coordinates, simply
a diagonal matrix: it stretches space along $n$ mutually perpendicular axes, each by a
real factor, with no rotation mixing the axes. We prove it in full, by induction, using
nothing beyond the previous chapter and the fact that a continuous function on a compact
set attains its maximum.
\end{intuition_box}

\section{Statement}

\begin{theorem}[Spectral Theorem for real symmetric matrices]\label{thm:spectral}
Let $A\in\R^{n\times n}$ be symmetric. Then there exist an orthogonal matrix
$Q=[\bq_1\ \cdots\ \bq_n]$ and real numbers $\lambda_1\ge\lambda_2\ge\cdots\ge\lambda_n$
such that
\[
A=Q\Lambda Q\T=\sum_{i=1}^n\lambda_i\,\bq_i\bq_i\T,\qquad
\Lambda=\operatorname{diag}(\lambda_1,\dots,\lambda_n),\qquad A\bq_i=\lambda_i\bq_i .
\]
That is, $A$ has an orthonormal basis of eigenvectors with real eigenvalues.
\end{theorem}

\section{A lemma: a maximizer of the Rayleigh quotient is an eigenvector}

\begin{lemma}[Existence of a real eigenpair]\label{lem:rayleigh}
Let $A\in\R^{n\times n}$ be symmetric. The function $\bx\mapsto\bx\T A\bx$ attains a
maximum on the unit sphere $S=\{\bx:\norm{\bx}=1\}$ at some $\bq\in S$, and that
maximizer satisfies $A\bq=\lambda\bq$ with $\lambda=\bq\T A\bq$.
\end{lemma}
\begin{proof}
The map $\bx\mapsto\bx\T A\bx$ is a polynomial, hence continuous, and $S$ is closed and
bounded (compact); therefore the maximum is attained at some $\bq\in S$, with value
$\lambda=\bq\T A\bq$. We show $\bq$ is an eigenvector by a first-order (Lagrange)
argument made elementary. Let $\bv$ be any vector with $\bv\perp\bq$, and consider the
unit-speed curve $\bx(t)=\cos(t)\,\bq+\sin(t)\,\hat\bv$ on $S$, where
$\hat\bv=\bv/\norm{\bv}$ (so $\norm{\bx(t)}=1$). Define $g(t)=\bx(t)\T A\bx(t)$. Since
$t=0$ maximizes $g$, $g'(0)=0$. Compute, using symmetry so that
$\tfrac{d}{dt}\bx\T A\bx=2\,\bx\T A\dot\bx$:
\[
g'(t)=2\,\bx(t)\T A\,\dot\bx(t),\qquad \dot\bx(0)=\hat\bv,\quad \bx(0)=\bq,
\]
so $0=g'(0)=2\,\bq\T A\hat\bv=2\inner{A\bq}{\hat\bv}$. Thus $A\bq\perp\bv$ for
\emph{every} $\bv\perp\bq$, which means $A\bq$ has no component orthogonal to $\bq$:
$A\bq\in\spn\{\bq\}$, say $A\bq=\lambda\bq$. Finally $\lambda=\bq\T A\bq$ by taking the
inner product with $\bq$.
\end{proof}

\section{Proof of the Spectral Theorem}

\begin{proof}[Proof of Theorem~\ref{thm:spectral}]
We induct on $n$. For $n=1$ every $1\times1$ matrix is diagonal and the claim is trivial
with $Q=[1]$.

Assume the theorem for dimension $n-1$. Given symmetric $A\in\R^{n\times n}$,
Lemma~\ref{lem:rayleigh} supplies a unit eigenvector $\bq_1$ with real eigenvalue
$\lambda_1=\bq_1\T A\bq_1$, chosen as the \emph{maximal} value of the Rayleigh quotient.
Let $W=\bq_1^\perp=\{\bx:\bq_1\T\bx=0\}$, an $(n-1)$-dimensional subspace.

\emph{$W$ is $A$-invariant.} If $\bx\in W$ then, using symmetry
(Lemma~\ref{lem:selfadj}),
\[
\inner{\bq_1}{A\bx}=\inner{A\bq_1}{\bx}=\lambda_1\inner{\bq_1}{\bx}=0,
\]
so $A\bx\in W$. Choose an orthonormal basis $\bw_2,\dots,\bw_n$ of $W$ and form the
orthogonal matrix $R=[\bq_1\ \bw_2\ \cdots\ \bw_n]$. Because $W$ is $A$-invariant and
$\bq_1$ is an eigenvector, the matrix of $A$ in this basis is block diagonal:
\[
R\T A R=\begin{bmatrix}\lambda_1 & \bzero\T\\ \bzero & A'\end{bmatrix},
\]
where $A'\in\R^{(n-1)\times(n-1)}$ represents $A|_W$. The off-diagonal blocks vanish:
the top row is $\bq_1\T A\bw_j=\lambda_1\bq_1\T\bw_j=0$, and the first column is
$\bw_j\T A\bq_1=\lambda_1\bw_j\T\bq_1=0$. Moreover $A'$ is symmetric, since
$R\T A R$ is symmetric ($A$ is) and a principal submatrix of a symmetric matrix is
symmetric.

By the inductive hypothesis $A'=Q'\Lambda'Q'^{\T}$ with $Q'$ orthogonal and $\Lambda'$
real diagonal. Set
\[
Q=R\begin{bmatrix}1 & \bzero\T\\ \bzero & Q'\end{bmatrix}.
\]
Then $Q$ is a product of orthogonal matrices, hence orthogonal
(Proposition~\ref{prop:Qprops}(iv)), and
\[
Q\T A Q=\begin{bmatrix}1&\bzero\T\\\bzero&Q'^{\T}\end{bmatrix}
(R\T A R)
\begin{bmatrix}1&\bzero\T\\\bzero&Q'\end{bmatrix}
=\begin{bmatrix}\lambda_1&\bzero\T\\\bzero&Q'^{\T}A'Q'\end{bmatrix}
=\begin{bmatrix}\lambda_1&\bzero\T\\\bzero&\Lambda'\end{bmatrix}=:\Lambda,
\]
a real diagonal matrix. Hence $A=Q\Lambda Q\T$. Reordering the columns of $Q$ and the
diagonal of $\Lambda$ so that $\lambda_1\ge\cdots\ge\lambda_n$ keeps $Q$ orthogonal and
completes the induction.
\end{proof}

\begin{corollary}[Orthogonal diagonalization and the variational characterization]
\label{cor:rayleigh}
With the ordering $\lambda_1\ge\cdots\ge\lambda_n$ of Theorem~\ref{thm:spectral},
\[
\lambda_1=\max_{\norm{\bx}=1}\bx\T A\bx,\qquad
\lambda_n=\min_{\norm{\bx}=1}\bx\T A\bx,
\]
and for any $\bx$, writing $\bx=\sum_i c_i\bq_i$ with $c_i=\bq_i\T\bx$,
\[
\bx\T A\bx=\sum_{i=1}^n\lambda_i c_i^2,\qquad \norm{\bx}^2=\sum_i c_i^2 .
\]
\end{corollary}
\begin{proof}
The expansions follow from $A\bq_i=\lambda_i\bq_i$ and orthonormality. The extremal
statements follow because $\sum_i\lambda_i c_i^2$ is a weighted average of the
$\lambda_i$ with weights $c_i^2\ge0$ summing to $\norm{\bx}^2=1$, maximized by putting
all weight on the largest $\lambda_1$ (achieved at $\bx=\bq_1$) and minimized at
$\bq_n$.
\end{proof}

\begin{geometry_box}
Theorem~\ref{thm:spectral} says: rotate space by $Q\T$ so that the eigenvectors become
the coordinate axes; in those coordinates $A$ is the diagonal matrix $\Lambda$, pure
axiswise scaling; then rotate back by $Q$. A symmetric matrix is a stretch along $n$
perpendicular axes --- nothing more.
\end{geometry_box}

\section*{Summary}
\addcontentsline{toc}{section}{Summary}
A maximizer of the Rayleigh quotient on the unit sphere is an eigenvector
(Lemma~\ref{lem:rayleigh}); deflating onto its orthogonal complement, which is
$A$-invariant, and inducting proves the Spectral Theorem
(Theorem~\ref{thm:spectral}): every real symmetric matrix is orthogonally diagonalizable
with real eigenvalues. Corollary~\ref{cor:rayleigh} records the variational
characterization that we will generalize (Ky Fan) in Chapter~\ref{ch:ey}.

\begin{exercise}
Prove the \emph{min--max} (Courant--Fischer) formula for $\lambda_k$ as a corollary of
the Spectral Theorem.
\end{exercise}
\begin{exercise}
Show $\tr(A)=\sum_i\lambda_i$ and $\det(A)=\prod_i\lambda_i$ for symmetric $A$.
\end{exercise}
\begin{exercise}
Diagonalize $A=\begin{psmallmatrix}2&1\\1&2\end{psmallmatrix}$ orthogonally and verify
$A=Q\Lambda Q\T$ explicitly.
\end{exercise}

%##################################################################
\chapter{Positive Semidefinite Matrices}\label{ch:psd}
%##################################################################

\begin{intuition_box}
The SVD is built from $A\T A$ and $AA\T$. These are not arbitrary symmetric matrices ---
they are \emph{positive semidefinite}, which forces their eigenvalues to be
\emph{nonnegative}. That nonnegativity is exactly what lets us take square roots and
define the singular values $\sigma_i=\sqrt{\lambda_i}$. This short chapter proves the
three facts the SVD construction will quote verbatim.
\end{intuition_box}

\begin{definition}[Positive semidefinite]
A symmetric matrix $M\in\R^{n\times n}$ is \emph{positive semidefinite} (PSD), written
$M\succeq0$, if $\bx\T M\bx\ge0$ for all $\bx\in\R^n$. It is \emph{positive definite}
($M\succ0$) if $\bx\T M\bx>0$ for all $\bx\ne\bzero$.
\end{definition}

\begin{proposition}[$A\T A$ and $AA\T$ are symmetric]\label{prop:gramsym}
For any $A\in\R^{m\times n}$, both $A\T A\in\R^{n\times n}$ and $AA\T\in\R^{m\times m}$
are symmetric.
\end{proposition}
\begin{proof}
$(A\T A)\T=A\T(A\T)\T=A\T A$, using $(XY)\T=Y\T X\T$ and $(A\T)\T=A$. Likewise
$(AA\T)\T=(A\T)\T A\T=AA\T$.
\end{proof}

\begin{proposition}[$A\T A$ and $AA\T$ are PSD]\label{prop:grampsd}
For any $A\in\R^{m\times n}$, $A\T A\succeq0$ and $AA\T\succeq0$. Hence (by the Spectral
Theorem) all their eigenvalues are $\ge0$.
\end{proposition}
\begin{proof}
For any $\bx\in\R^n$, $\bx\T(A\T A)\bx=(A\bx)\T(A\bx)=\norm{A\bx}^2\ge0$, so
$A\T A\succeq0$. Replacing $A$ by $A\T$ gives $AA\T\succeq0$. If $M\bq=\lambda\bq$ with
$\norm{\bq}=1$ and $M\succeq0$, then $\lambda=\bq\T M\bq\ge0$.
\end{proof}

\begin{proposition}[Shared nonzero spectrum of $A\T A$ and $AA\T$]\label{prop:sharedspec}
The matrices $A\T A$ and $AA\T$ have the same nonzero eigenvalues, with the same
multiplicities. Specifically, if $A\T A\,\bv=\lambda\bv$ with $\lambda\ne0$, then
$\bu:=A\bv$ satisfies $AA\T\bu=\lambda\bu$ and $\bu\ne\bzero$.
\end{proposition}
\begin{proof}
From $A\T A\bv=\lambda\bv$, left-multiply by $A$:
$AA\T(A\bv)=A(A\T A\bv)=A(\lambda\bv)=\lambda(A\bv)$, i.e.\ $AA\T\bu=\lambda\bu$.
Moreover $\norm{A\bv}^2=\bv\T A\T A\bv=\lambda\norm{\bv}^2\ne0$, so $\bu=A\bv\ne\bzero$.
The map $\bv\mapsto A\bv$ thus carries each $\lambda$-eigenvector of $A\T A$ to a
$\lambda$-eigenvector of $AA\T$; it is injective on that eigenspace (if $A\bv=\bzero$
then $\lambda\bv=A\T A\bv=\bzero$ forces $\bv=\bzero$), so multiplicities match. By
symmetry of roles, the correspondence is a bijection on nonzero eigenspaces.
\end{proof}

\begin{ml_box}
Proposition~\ref{prop:grampsd} is the reason PCA works at all: the covariance matrix
$\Sigma=\tfrac1m X_c\T X_c$ is of the form $A\T A$, hence symmetric PSD, hence has real
nonnegative eigenvalues $\lambda_i\ge0$ that we interpret as \emph{variances}. A negative
variance would be nonsense; PSD-ness guarantees it never happens.
Proposition~\ref{prop:sharedspec} is the engine of the ``Gram-matrix trick'': when there
are far fewer samples than features, one eigendecomposes the small $m\times m$ matrix
$AA\T$ and transfers the eigenvectors to the large space via $\bv\mapsto A\bv$.
\end{ml_box}

\section*{Summary}
\addcontentsline{toc}{section}{Summary}
$A\T A$ and $AA\T$ are always symmetric (Proposition~\ref{prop:gramsym}) and PSD
(Proposition~\ref{prop:grampsd}), so their eigenvalues are real and nonnegative ---
permitting the square roots $\sigma_i=\sqrt{\lambda_i}$ that define singular values. The
two products share their nonzero spectrum (Proposition~\ref{prop:sharedspec}), the link
that will tie left and right singular vectors together.

\begin{exercise}
Prove $A\T A\succ0$ (positive \emph{definite}) iff $A$ has linearly independent columns
(equivalently $\nul(A)=\{\bzero\}$).
\end{exercise}
\begin{exercise}
Show that $M\succeq0$ iff $M=B\T B$ for some matrix $B$. \emph{Hint:} use the Spectral
Theorem to build $B=\Lambda^{1/2}Q\T$.
\end{exercise}
\begin{exercise}
For $A=\begin{psmallmatrix}3&0\\4&5\end{psmallmatrix}$, compute $A\T A$ and $AA\T$ and
verify they share the eigenvalues $\{45,5\}$.
\end{exercise}

%##################################################################
\chapter{Derivation of the Singular Value Decomposition}\label{ch:svd}
%##################################################################

\begin{intuition_box}
We now \emph{build} the SVD; we do not assume it. The plan: $A\T A$ is symmetric PSD
(Chapter~\ref{ch:psd}), so the Spectral Theorem (Chapter~\ref{ch:spectral}) gives it an
orthonormal eigenbasis with nonnegative eigenvalues. We declare those eigenvectors to be
the right singular vectors, the square roots of the eigenvalues to be the singular
values, and we \emph{derive} the left singular vectors as the normalized images
$A\bv_i/\sigma_i$. Verifying orthonormality of the left vectors and reassembling proves
$A=U\Sigma V\T$. Because the construction only ever uses the Spectral Theorem applied to
a matrix that always exists, the SVD always exists.
\end{intuition_box}

\section{Construction}

\begin{theorem}[Existence of the SVD]\label{thm:svdexist}
Every $A\in\R^{m\times n}$ admits a factorization
\[
A=U\Sigma V\T,
\]
where $U\in\R^{m\times m}$ and $V\in\R^{n\times n}$ are orthogonal and
$\Sigma\in\R^{m\times n}$ has entries $\Sigma_{ii}=\sigma_i\ge0$ for
$i=1,\dots,\min(m,n)$ and zeros off the diagonal, with
$\sigma_1\ge\sigma_2\ge\cdots\ge\sigma_r>0=\sigma_{r+1}=\cdots$, where $r=\rank(A)$.
\end{theorem}
\begin{proof}
\textbf{Right singular vectors and singular values.}
By Propositions~\ref{prop:gramsym}--\ref{prop:grampsd}, $A\T A\in\R^{n\times n}$ is
symmetric PSD. The Spectral Theorem (Theorem~\ref{thm:spectral}) gives an orthonormal
basis $\bv_1,\dots,\bv_n$ of $\R^n$ with
\[
A\T A\,\bv_i=\lambda_i\bv_i,\qquad \lambda_1\ge\cdots\ge\lambda_n\ge0 .
\]
Define $\sigma_i=\sqrt{\lambda_i}\ge0$. Let $r$ be the number of strictly positive
$\sigma_i$ (we show $r=\rank A$ below). Set $V=[\bv_1\ \cdots\ \bv_n]$, orthogonal.

\textbf{Left singular vectors.}
For $i=1,\dots,r$ (so $\sigma_i>0$) define
\[
\bu_i:=\frac{1}{\sigma_i}A\bv_i\in\R^m .
\]
These are orthonormal: for $i,j\le r$,
\[
\bu_i\T\bu_j=\frac{1}{\sigma_i\sigma_j}(A\bv_i)\T(A\bv_j)
=\frac{1}{\sigma_i\sigma_j}\bv_i\T(A\T A\bv_j)
=\frac{\lambda_j}{\sigma_i\sigma_j}\bv_i\T\bv_j
=\frac{\sigma_j^2}{\sigma_i\sigma_j}\delta_{ij}
=\frac{\sigma_j}{\sigma_i}\delta_{ij}=\delta_{ij}.
\]
Extend $\bu_1,\dots,\bu_r$ to an orthonormal basis $\bu_1,\dots,\bu_m$ of $\R^m$
(Gram--Schmidt on any completion), and set $U=[\bu_1\ \cdots\ \bu_m]$, orthogonal.

\textbf{The annihilated directions.}
For $i>r$ we have $\lambda_i=0$, so
$\norm{A\bv_i}^2=\bv_i\T A\T A\bv_i=\lambda_i=0$, hence $A\bv_i=\bzero$.

\textbf{Reassembly.}
Define $\Sigma\in\R^{m\times n}$ by $\Sigma_{ii}=\sigma_i$ and zeros elsewhere. We verify
$A=U\Sigma V\T$ by checking both sides agree on the basis $\bv_1,\dots,\bv_n$ (two linear
maps equal on a basis are equal, Lemma~\ref{lem:uniquecoord}). The right side sends
\[
(U\Sigma V\T)\bv_j=U\Sigma\be_j=U(\sigma_j\be_j)=\sigma_j\bu_j\quad(j\le r),\qquad
(U\Sigma V\T)\bv_j=\sigma_j\bu_j=\bzero\quad(j>r,\ \sigma_j=0),
\]
using $V\T\bv_j=\be_j$. The left side sends $A\bv_j=\sigma_j\bu_j$ for $j\le r$ (by
definition of $\bu_j$) and $A\bv_j=\bzero$ for $j>r$. The two maps agree on the basis,
so $A=U\Sigma V\T$.

\textbf{$r=\rank A$.}
From $A=U\Sigma V\T$ with $U,V$ invertible, $\rank(A)=\rank(\Sigma)$ (multiplying by
invertible matrices preserves rank, Exercise), and $\Sigma$ has exactly $r$ nonzero
entries on its diagonal, so $\rank(\Sigma)=r$.
\end{proof}

\begin{corollary}[Outer-product / dyadic form]\label{cor:dyadic}
With the notation of Theorem~\ref{thm:svdexist},
\[
A=\sum_{i=1}^{r}\sigma_i\,\bu_i\bv_i\T .
\]
Each $\bu_i\bv_i\T$ is rank one (Theorem~\ref{thm:rank1}) and the terms are ordered by
$\sigma_1\ge\cdots\ge\sigma_r>0$.
\end{corollary}
\begin{proof}
Multiplying out $U\Sigma V\T$ keeps only the $r$ nonzero diagonal entries of $\Sigma$,
each pairing column $\bu_i$ of $U$ with row $\bv_i\T$ of $V\T$, weighted by $\sigma_i$.
\end{proof}

\begin{corollary}[$U$ and $V$ as eigenvectors; the EVD link]\label{cor:evdlink}
The columns of $V$ are orthonormal eigenvectors of $A\T A$ with eigenvalues $\sigma_i^2$,
and the columns of $U$ are orthonormal eigenvectors of $AA\T$ with the same nonzero
eigenvalues:
\[
A\T A=V\Sigma\T\Sigma V\T,\qquad AA\T=U\Sigma\Sigma\T U\T .
\]
\end{corollary}
\begin{proof}
$A\T A=(U\Sigma V\T)\T(U\Sigma V\T)=V\Sigma\T U\T U\Sigma V\T=V(\Sigma\T\Sigma)V\T$,
using $U\T U=I$. This is the orthogonal diagonalization of $A\T A$, so the columns of $V$
are its eigenvectors with eigenvalues $(\Sigma\T\Sigma)_{ii}=\sigma_i^2$. Symmetrically
$AA\T=U(\Sigma\Sigma\T)U\T$.
\end{proof}

\section{Full, reduced, and compact SVD}

\begin{definition}[Three forms]\label{def:svdforms}
Let $r=\rank A$.
\begin{itemize}[leftmargin=1.6em]
\item \textbf{Full SVD:} $A=U\Sigma V\T$ with $U\in\R^{m\times m}$,
$V\in\R^{n\times n}$ orthogonal and $\Sigma\in\R^{m\times n}$.
\item \textbf{Reduced (thin) SVD:} keep the first $\min(m,n)$ columns of $U$ and the
corresponding block of $\Sigma$.
\item \textbf{Compact SVD:} keep only the $r$ columns with $\sigma_i>0$:
$A=U_r\Sigma_r V_r\T$, $U_r\in\R^{m\times r}$, $\Sigma_r\in\R^{r\times r}$
(invertible diagonal), $V_r\in\R^{n\times r}$, $U_r\T U_r=I_r=V_r\T V_r$.
\end{itemize}
\end{definition}

\begin{proposition}[The four fundamental subspaces from the SVD]\label{prop:foursub}
With the compact SVD $A=U_r\Sigma_r V_r\T$,
\[
\col(A)=\spn\{\bu_1,\dots,\bu_r\},\qquad
\nul(A\T)=\spn\{\bu_{r+1},\dots,\bu_m\},
\]
\[
\row(A)=\spn\{\bv_1,\dots,\bv_r\},\qquad
\nul(A)=\spn\{\bv_{r+1},\dots,\bv_n\}.
\]
\end{proposition}
\begin{proof}
$A\bx=\sum_{i\le r}\sigma_i(\bv_i\T\bx)\bu_i$ lies in $\spn\{\bu_i\}_{i\le r}$, and every
$\bu_i$ ($i\le r$) is attained (take $\bx=\bv_i$), so $\col(A)=\spn\{\bu_i\}_{i\le r}$.
For the null space, $A\bx=\bzero$ iff $\sigma_i(\bv_i\T\bx)=0$ for all $i\le r$ iff
$\bv_i\T\bx=0$ for $i\le r$ (as $\sigma_i>0$), i.e.\ $\bx\in\spn\{\bv_i\}_{i>r}$. The
$A\T$ statements follow by applying these to $A\T=V\Sigma\T U\T$.
\end{proof}

\section{Uniqueness}

\begin{theorem}[Uniqueness of the SVD]\label{thm:svduniq}
The singular values $\sigma_1\ge\cdots\ge\sigma_r>0$ are uniquely determined by $A$
(they are the square roots of the nonzero eigenvalues of $A\T A$). If the nonzero
singular values are \emph{distinct}, then each pair $(\bu_i,\bv_i)$ is unique up to a
common sign $(\bu_i,\bv_i)\mapsto(-\bu_i,-\bv_i)$. If a singular value is repeated, the
corresponding singular vectors are determined only up to an orthonormal rotation within
the shared eigenspace.
\end{theorem}
\begin{proof}
The $\sigma_i^2$ are the eigenvalues of the matrix $A\T A$, which depends only on $A$;
eigenvalues are unique, so the $\sigma_i$ are. For a simple (multiplicity-one) eigenvalue
$\sigma_i^2$ of $A\T A$, the eigenspace is one-dimensional, so $\bv_i$ is determined up
to sign; then $\bu_i=A\bv_i/\sigma_i$ inherits the same sign choice, giving the common
sign freedom. For a repeated eigenvalue, the Spectral Theorem only fixes the eigenspace,
within which any orthonormal basis is admissible; the relation $\bu_i=A\bv_i/\sigma_i$
transports that freedom to the $\bu_i$.
\end{proof}

\begin{example}[A complete $2\times2$ SVD by hand]\label{ex:svd22}
Let $A=\begin{psmallmatrix}3&0\\4&5\end{psmallmatrix}$.
Then $A\T A=\begin{psmallmatrix}25&20\\20&25\end{psmallmatrix}$, with characteristic
equation $(25-\lambda)^2-400=0$, so $\lambda=25\pm20$, giving $\lambda_1=45$,
$\lambda_2=5$ and $\sigma_1=3\sqrt5$, $\sigma_2=\sqrt5$. The eigenvectors of $A\T A$ are
$\bv_1=\tfrac1{\sqrt2}(1,1)\T$ (for $45$) and $\bv_2=\tfrac1{\sqrt2}(1,-1)\T$ (for $5$).
Then
\[
\bu_1=\frac{A\bv_1}{\sigma_1}=\frac{1}{3\sqrt5}\cdot\frac1{\sqrt2}
\begin{psmallmatrix}3\\9\end{psmallmatrix}=\frac1{\sqrt{10}}
\begin{psmallmatrix}1\\3\end{psmallmatrix},\qquad
\bu_2=\frac{A\bv_2}{\sigma_2}=\frac{1}{\sqrt5}\cdot\frac1{\sqrt2}
\begin{psmallmatrix}3\\-1\end{psmallmatrix}=\frac1{\sqrt{10}}
\begin{psmallmatrix}3\\-1\end{psmallmatrix}.
\]
These $\bu_i$ are exactly the eigenvectors of $AA\T=\begin{psmallmatrix}9&12\\12&41
\end{psmallmatrix}$ (eigenvalues $45,5$), confirming Corollary~\ref{cor:evdlink}. One
checks $A=\sigma_1\bu_1\bv_1\T+\sigma_2\bu_2\bv_2\T$ reproduces $A$.
\end{example}

\begin{geometry_box}
Reading $A=U\Sigma V\T$ right to left: $V\T$ \emph{rotates} the input so the right
singular vectors become coordinate axes; $\Sigma$ \emph{scales} axis $i$ by $\sigma_i$
(and embeds/projects between $\R^n$ and $\R^m$); $U$ \emph{rotates} the result into the
output space along the left singular vectors. Consequently $A$ maps the unit sphere of
$\R^n$ onto an ellipsoid in $\R^m$ whose semi-axes have lengths $\sigma_i$ in directions
$\bu_i$.
\end{geometry_box}

\begin{miscon_box}
``The SVD is just the eigendecomposition of $A$.'' No: $A$ need not be square or
diagonalizable, and even when it is, the singular values are \emph{not} the absolute
values of the eigenvalues in general (they coincide only for symmetric PSD $A$). The SVD
is the eigendecomposition of the symmetric PSD matrices $A\T A$ and $AA\T$, stitched
together --- a strictly more general and always-available object.
\end{miscon_box}

\section*{Summary}
\addcontentsline{toc}{section}{Summary}
Applying the Spectral Theorem to $A\T A$ yields $V$ and the $\sigma_i$; defining
$\bu_i=A\bv_i/\sigma_i$ yields an orthonormal $U$; reassembly gives $A=U\Sigma V\T$
(Theorem~\ref{thm:svdexist}), equivalently the dyadic sum (Corollary~\ref{cor:dyadic}).
$V,U$ diagonalize $A\T A,AA\T$ (Corollary~\ref{cor:evdlink}); the SVD reads off all four
fundamental subspaces (Proposition~\ref{prop:foursub}); it is unique up to the sign /
rotation freedom of Theorem~\ref{thm:svduniq}.

\begin{exercise}
Prove that multiplying by invertible (in particular orthogonal) matrices preserves rank,
the fact quoted at the end of Theorem~\ref{thm:svdexist}.
\end{exercise}
\begin{exercise}
Show the pseudoinverse $A^+:=V_r\Sigma_r\inv U_r\T$ satisfies $AA^+A=A$ and $A^+AA^+=A^+$.
\end{exercise}
\begin{exercise}
For $A$ symmetric PSD, prove the SVD and the EVD coincide ($U=V=Q$, $\sigma_i=\lambda_i$).
\end{exercise}
\begin{exercise}
Compute the full SVD of $A=\begin{psmallmatrix}1&1&0\\0&1&1\end{psmallmatrix}$.
\end{exercise}

%##################################################################
\chapter{Matrix Norms}\label{ch:norms}
%##################################################################

\begin{intuition_box}
To say one matrix is the \emph{best} approximation of another we must measure the size of
their difference. Two norms matter. The \emph{Frobenius} norm treats a matrix as a long
vector and measures total entrywise energy; it is the natural least-squares error. The
\emph{spectral} norm measures the largest stretch $A$ applies to any unit vector; it is
the natural worst-case error. The Eckart--Young--Mirsky theorem holds for \emph{both},
and we prove the identities we need for each.
\end{intuition_box}

\section{The Frobenius norm and trace}

\begin{definition}[Trace]
For square $M$, $\tr(M)=\sum_i M_{ii}$.
\end{definition}

\begin{lemma}[Cyclic invariance of trace]\label{lem:cyc}
For conformable matrices, $\tr(XY)=\tr(YX)$, and hence $\tr(XYZ)=\tr(YZX)=\tr(ZXY)$.
\end{lemma}
\begin{proof}
$\tr(XY)=\sum_i\sum_j X_{ij}Y_{ji}=\sum_j\sum_i Y_{ji}X_{ij}=\tr(YX)$. The three-factor
identity follows by grouping $X(YZ)$ etc.
\end{proof}

\begin{definition}[Frobenius norm]
$\displaystyle \fn{A}=\Bigl(\sum_{i=1}^m\sum_{j=1}^n a_{ij}^2\Bigr)^{1/2}$ for
$A\in\R^{m\times n}$.
\end{definition}

\begin{proposition}[Frobenius via trace]\label{prop:frobtr}
$\fn{A}^2=\tr(A\T A)=\tr(AA\T)=\sum_j\norm{\ba_j}^2$, where $\ba_j$ are the columns.
\end{proposition}
\begin{proof}
$(A\T A)_{jj}=\sum_i a_{ij}^2=\norm{\ba_j}^2$; summing over $j$ gives
$\tr(A\T A)=\sum_{i,j}a_{ij}^2=\fn{A}^2$. The equality $\tr(A\T A)=\tr(AA\T)$ is
Lemma~\ref{lem:cyc}.
\end{proof}

\begin{theorem}[Orthogonal invariance of the Frobenius norm]\label{thm:frobinv}
If $Q\in\R^{m\times m}$, $P\in\R^{n\times n}$ are orthogonal, then for all
$A\in\R^{m\times n}$, $\ \fn{QAP}=\fn{A}$.
\end{theorem}
\begin{proof}
Using Proposition~\ref{prop:frobtr} and the cyclic property,
\[
\fn{QAP}^2=\tr\bigl((QAP)\T QAP\bigr)=\tr\bigl(P\T A\T Q\T QAP\bigr)
=\tr\bigl(A\T A\,PP\T\bigr)=\tr(A\T A)=\fn{A}^2,
\]
where $Q\T Q=I$ and $PP\T=I$.
\end{proof}

\begin{corollary}[Frobenius norm in singular values]\label{cor:frobsing}
$\displaystyle \fn{A}^2=\sum_{i=1}^r\sigma_i^2$.
\end{corollary}
\begin{proof}
By Theorem~\ref{thm:frobinv} with the SVD $A=U\Sigma V\T$,
$\fn{A}=\fn{U\Sigma V\T}=\fn{\Sigma}$, and $\fn{\Sigma}^2=\sum_i\sigma_i^2$.
\end{proof}

\section{The spectral (operator-2) norm}

\begin{definition}[Spectral norm]
$\displaystyle \sn{A}=\max_{\bx\ne\bzero}\frac{\norm{A\bx}}{\norm{\bx}}
=\max_{\norm{\bx}=1}\norm{A\bx}$.
\end{definition}

\begin{theorem}[Spectral norm equals the top singular value]\label{thm:specnorm}
$\sn{A}=\sigma_1$, the largest singular value, with the maximum attained at
$\bx=\bv_1$ (and $A\bv_1=\sigma_1\bu_1$).
\end{theorem}
\begin{proof}
For $\norm{\bx}=1$ write $\bx=\sum_i c_i\bv_i$ with $\sum_i c_i^2=1$
(Corollary~\ref{cor:rayleigh} applied to the orthonormal basis $V$). Then
$A\bx=\sum_{i\le r}\sigma_i c_i\bu_i$, and since the $\bu_i$ are orthonormal,
\[
\norm{A\bx}^2=\sum_{i\le r}\sigma_i^2 c_i^2\le\sigma_1^2\sum_i c_i^2=\sigma_1^2 ,
\]
so $\norm{A\bx}\le\sigma_1$. Equality holds at $\bx=\bv_1$ (then $c_1=1$, rest $0$),
where $\norm{A\bv_1}=\sigma_1$. Hence the maximum is exactly $\sigma_1$.
\end{proof}

\begin{remark}
Both norms are written entirely in singular values:
$\fn{A}^2=\sum_i\sigma_i^2$ and $\sn{A}=\sigma_1$. This is why the SVD is the right tool
for approximation questions phrased in either norm.
\end{remark}

\section*{Summary}
\addcontentsline{toc}{section}{Summary}
The Frobenius norm is total entrywise energy, equal to $\tr(A\T A)$
(Proposition~\ref{prop:frobtr}), invariant under orthogonal multiplication
(Theorem~\ref{thm:frobinv}), and equal to $\sqrt{\sum_i\sigma_i^2}$
(Corollary~\ref{cor:frobsing}). The spectral norm is the maximal stretch, equal to
$\sigma_1$ (Theorem~\ref{thm:specnorm}). Both are functions of the singular values
alone.

\begin{exercise}
Prove $\sn{A}\le\fn{A}\le\sqrt{r}\,\sn{A}$ where $r=\rank A$.
\end{exercise}
\begin{exercise}
Show $\fn{AB}\le\sn{A}\fn{B}$ and $\fn{AB}\le\fn{A}\sn{B}$.
\end{exercise}
\begin{exercise}
Prove $\sn{A}=\sn{A\T}$ and $\fn{A}=\fn{A\T}$ directly from the SVD.
\end{exercise}

%##################################################################
\chapter{The Best Rank-$k$ Approximation: Eckart--Young--Mirsky}\label{ch:ey}
%##################################################################

\begin{intuition_box}
Everything has been building to this. We truncate the dyadic SVD to its first $k$ terms
and prove that the result is the closest rank-$k$ matrix to $A$ --- not merely a good
choice, but the provably optimal one --- in both the Frobenius and the spectral norm. The
Frobenius proof goes through two lemmas: ``best matrix for a fixed column space is a
projection'' (which converts the problem from choosing a matrix to choosing a subspace),
and the \emph{Ky Fan} maximum principle (which solves the subspace problem). The spectral
proof uses a short dimension-counting argument. No step is taken on faith.
\end{intuition_box}

\section{The truncated SVD and its error}

\begin{definition}[Truncated SVD]
For $A=\sum_{i=1}^r\sigma_i\bu_i\bv_i\T$ and $1\le k\le r$, define
\[
A_k=\sum_{i=1}^k\sigma_i\bu_i\bv_i\T=U_k\Sigma_k V_k\T,
\]
where $U_k,V_k$ keep the first $k$ columns and $\Sigma_k=\operatorname{diag}
(\sigma_1,\dots,\sigma_k)$. By construction $\rank(A_k)\le k$.
\end{definition}

\begin{proposition}[Error of the truncation]\label{prop:truncerr}
$\displaystyle \fn{A-A_k}^2=\sum_{i=k+1}^{r}\sigma_i^2$ and
$\sn{A-A_k}=\sigma_{k+1}$ (with $\sigma_{r+1}:=0$).
\end{proposition}
\begin{proof}
$A-A_k=\sum_{i=k+1}^r\sigma_i\bu_i\bv_i\T$ is itself an SVD (orthonormal $\bu_i,\bv_i$)
with singular values $\sigma_{k+1}\ge\cdots\ge\sigma_r$. Apply
Corollary~\ref{cor:frobsing} for the Frobenius statement and
Theorem~\ref{thm:specnorm} for the spectral statement.
\end{proof}

This computes the error of \emph{our} candidate. The theorem asserts no rank-$k$ matrix
does better. We prove that next.

\section{Lemma A: reduction to a subspace problem}

\begin{lemma}[Best matrix for a fixed column space]\label{lem:projbest}
Fix a subspace $S\subseteq\R^m$ with $\dimm S\le k$, and let $P_S$ be the orthogonal
projector onto $S$. Among all $B\in\R^{m\times n}$ with $\col(B)\subseteq S$, the
Frobenius error $\fn{A-B}$ is minimized by $B=P_SA$, and
\[
\min_{\col(B)\subseteq S}\fn{A-B}^2=\fn{(I-P_S)A}^2=\fn{A}^2-\fn{P_SA}^2 .
\]
\end{lemma}
\begin{proof}
Every column $\bb_j$ of $B$ lies in $S$. By the best-approximation theorem
(Theorem~\ref{thm:bestapprox}(c)) applied to the column $\ba_j$ of $A$,
$\norm{\ba_j-\bb_j}\ge\norm{\ba_j-P_S\ba_j}$ with equality at $\bb_j=P_S\ba_j$.
Squaring and summing over $j$ (Proposition~\ref{prop:frobtr}) gives
$\fn{A-B}^2\ge\fn{(I-P_S)A}^2$, attained at $B=P_SA$. The final equality is the matrix
Pythagoras identity (Corollary~\ref{cor:matpyth}).
\end{proof}

\begin{remark}[Why this is the crucial reduction]
A matrix has rank $\le k$ \emph{iff} its columns lie in some subspace of dimension
$\le k$ (Theorem~\ref{thm:rankequiv}). Lemma~\ref{lem:projbest} says that once the
subspace $S$ is fixed, the optimal $B$ is forced. Hence
\[
\min_{\rank(B)\le k}\fn{A-B}^2
=\min_{\dimm S\le k}\bigl(\fn{A}^2-\fn{P_SA}^2\bigr)
=\fn{A}^2-\max_{\dimm S\le k}\fn{P_SA}^2 .
\]
Minimizing the error is the \emph{same} as maximizing the captured energy
$\fn{P_SA}^2$ over $k$-dimensional subspaces. That maximization is Lemma~B.
\end{remark}

\section{Lemma B: the Ky Fan maximum principle}

We first rewrite the captured energy as a sum of quadratic forms.

\begin{lemma}[Captured energy is a sum of quadratic forms]\label{lem:capture}
If $S$ has orthonormal basis $\bw_1,\dots,\bw_d$ (so $P_S=\sum_j\bw_j\bw_j\T$), then
\[
\fn{P_SA}^2=\sum_{j=1}^d\bw_j\T(AA\T)\bw_j .
\]
\end{lemma}
\begin{proof}
$\fn{P_SA}^2=\tr\bigl((P_SA)\T P_SA\bigr)=\tr(A\T P_S\T P_S A)=\tr(A\T P_S A)$, using
$P_S\T=P_S$ and $P_S^2=P_S$ (Proposition~\ref{prop:projprops}). With
$P_S=\sum_j\bw_j\bw_j\T$ and cyclicity of trace,
$\tr(A\T P_S A)=\sum_j\tr(A\T\bw_j\bw_j\T A)=\sum_j(\bw_j\T A)(A\T\bw_j)
=\sum_j\bw_j\T(AA\T)\bw_j$.
\end{proof}

\begin{lemma}[Ky Fan: maximal sum of quadratic forms]\label{lem:kyfan}
Let $M\in\R^{m\times m}$ be symmetric with eigenvalues
$\mu_1\ge\mu_2\ge\cdots\ge\mu_m$ and orthonormal eigenvectors $\bq_1,\dots,\bq_m$
(Spectral Theorem). For every orthonormal set $\bw_1,\dots,\bw_k$,
\[
\sum_{j=1}^k\bw_j\T M\bw_j\ \le\ \mu_1+\cdots+\mu_k,
\]
with equality when $\bw_j=\bq_j$.
\end{lemma}
\begin{proof}
Extend $\bw_1,\dots,\bw_k$ to an orthonormal basis $\bw_1,\dots,\bw_m$ of $\R^m$. Set
$c_{ij}=\bq_i\T\bw_j$. Using $M=\sum_i\mu_i\bq_i\bq_i\T$,
\[
\sum_{j=1}^k\bw_j\T M\bw_j
=\sum_{j=1}^k\sum_{i=1}^m\mu_i(\bq_i\T\bw_j)^2
=\sum_{i=1}^m\mu_i\,t_i,\qquad t_i:=\sum_{j=1}^k c_{ij}^2 .
\]
We claim $0\le t_i\le1$ and $\sum_i t_i=k$:
\begin{itemize}[leftmargin=1.5em]
\item $t_i\ge0$ trivially; and $t_i\le\sum_{j=1}^m c_{ij}^2=\norm{\bq_i}^2=1$, because
$\{\bw_j\}_{j=1}^m$ is a full orthonormal basis so the squared coordinates of the unit
vector $\bq_i$ sum to $1$; truncating to $k\le m$ terms can only decrease the sum.
\item $\sum_{i=1}^m t_i=\sum_{j=1}^k\sum_{i=1}^m c_{ij}^2=\sum_{j=1}^k\norm{\bw_j}^2=k$,
since $\{\bq_i\}$ is an orthonormal basis.
\end{itemize}
We must therefore bound $\sum_i\mu_it_i$ over $0\le t_i\le1,\ \sum_i t_i=k$. For any such
$t$,
\[
\sum_{i=1}^m\mu_it_i-\sum_{i=1}^k\mu_i
=\sum_{i=1}^k\mu_i(t_i-1)+\sum_{i>k}\mu_it_i
\le\mu_k\sum_{i=1}^k(t_i-1)+\mu_k\sum_{i>k}t_i
=\mu_k\Bigl(\sum_i t_i-k\Bigr)=0,
\]
where the inequality uses $\mu_i\ge\mu_k$ together with $t_i-1\le0$ on the first sum, and
$\mu_i\le\mu_k$ with $t_i\ge0$ on the second. Hence
$\sum_i\mu_it_i\le\sum_{i=1}^k\mu_i$. Choosing $\bw_j=\bq_j$ gives $t_i=1$ for $i\le k$
and $0$ otherwise, achieving equality.
\end{proof}

\begin{geometry_box}
Ky Fan in words: each unit direction $\bw_j$ has a total ``budget'' of $1$ to distribute
across the eigen-axes of $M$; $k$ orthonormal directions distribute a total budget of
$k$, but no single axis can absorb more than $1$. To maximize $\sum_i\mu_it_i$ you place
a full unit on each of the $k$ largest eigenvalues --- i.e.\ align your directions with
the top $k$ eigenvectors.
\end{geometry_box}

\section{The theorem (Frobenius norm)}

\begin{theorem}[Eckart--Young--Mirsky, Frobenius norm]\label{thm:eyf}
Let $A\in\R^{m\times n}$ have singular values $\sigma_1\ge\cdots\ge\sigma_r>0$ and let
$1\le k\le r$. Then for every $B$ with $\rank(B)\le k$,
\[
\fn{A-B}\ \ge\ \fn{A-A_k},\qquad\text{i.e.}\qquad
\min_{\rank(B)\le k}\fn{A-B}^2=\sum_{i=k+1}^{r}\sigma_i^2 ,
\]
attained at $B=A_k$. If $\sigma_k>\sigma_{k+1}$ the minimizer $A_k$ is unique.
\end{theorem}
\begin{proof}
Let $B$ have $\rank(B)\le k$ and set $S=\col(B)$, so $\dimm S\le k$. By
Lemma~\ref{lem:projbest},
\[
\fn{A-B}^2\ \ge\ \fn{A}^2-\fn{P_SA}^2 .
\]
Enlarging $S$ to dimension exactly $k$ only increases $\fn{P_SA}^2$
(Lemma~\ref{lem:capture} adds nonnegative terms), so assume $\dimm S=k$ with orthonormal
basis $\bw_1,\dots,\bw_k$. By Lemma~\ref{lem:capture} and the Ky Fan principle
(Lemma~\ref{lem:kyfan}) applied to $M=AA\T$ --- whose eigenvalues are
$\sigma_1^2\ge\cdots\ge\sigma_r^2\ge0$ with eigenvectors the left singular vectors $\bu_i$
(Corollary~\ref{cor:evdlink}) ---
\[
\fn{P_SA}^2=\sum_{j=1}^k\bw_j\T(AA\T)\bw_j\ \le\ \sigma_1^2+\cdots+\sigma_k^2 .
\]
Therefore, using $\fn{A}^2=\sum_{i=1}^r\sigma_i^2$ (Corollary~\ref{cor:frobsing}),
\[
\fn{A-B}^2\ \ge\ \sum_{i=1}^r\sigma_i^2-\sum_{i=1}^k\sigma_i^2=\sum_{i=k+1}^r\sigma_i^2
=\fn{A-A_k}^2 ,
\]
the last equality by Proposition~\ref{prop:truncerr}. Thus no rank-$k$ matrix beats
$A_k$. The bound is attained by $A_k$: taking
$S^\star=\spn\{\bu_1,\dots,\bu_k\}$ makes Ky Fan an equality, and the optimal $B$ from
Lemma~\ref{lem:projbest} is
\[
P_{S^\star}A=U_kU_k\T A=U_kU_k\T(U\Sigma V\T)=U_k\Sigma_k V_k\T=A_k,
\]
because $U_k\T U$ extracts the first $k$ rows of $\Sigma V\T$.

\emph{Uniqueness when $\sigma_k>\sigma_{k+1}$.} Equality throughout forces (i) Ky Fan
tight and (ii) the projection bound tight. With the strict gap
$\sigma_k^2>\sigma_{k+1}^2$, the optimal weight vector $t$ in Lemma~\ref{lem:kyfan} is the
unique vertex $t_1=\cdots=t_k=1$, which forces $S=\spn\{\bu_1,\dots,\bu_k\}$ exactly;
then (ii) forces $B=P_SA=A_k$.
\end{proof}

\begin{miscon_box}
A popular shortcut argues: rotate by $U,V$ to reduce the problem to
$\min_{\rank(C)\le k}\fn{\Sigma-C}$ with $C=U\T BV$; observe that off-diagonal entries of
$C$ only add error; conclude ``the optimal $C$ is diagonal,'' then keep the top $k$
$\sigma_i$. The \emph{answer} is right but the conclusion ``optimal $C$ is diagonal'' is
\emph{not} justified by the off-diagonal observation alone: zeroing off-diagonals can
change the rank, so one is not comparing matrices of the same rank. Optimizing the
diagonal and the off-diagonal separately ignores that the rank constraint couples all
entries. The genuine content --- that no rank-$k$ competitor (diagonal or not) beats the
truncation --- is supplied precisely by Ky Fan (Lemma~\ref{lem:kyfan}). Our proof never
assumes diagonality; the diagonal form of $A_k$ \emph{emerges} as the conclusion.
\end{miscon_box}

\section{The theorem (spectral norm)}

\begin{theorem}[Eckart--Young--Mirsky, spectral norm]\label{thm:eys}
With the notation above, for every $B$ with $\rank(B)\le k$,
\[
\sn{A-B}\ \ge\ \sigma_{k+1}=\sn{A-A_k}.
\]
Hence $A_k$ is also a best rank-$k$ approximation in the spectral norm.
\end{theorem}
\begin{proof}
We already have $\sn{A-A_k}=\sigma_{k+1}$ (Proposition~\ref{prop:truncerr}). Let $B$ have
$\rank(B)\le k$, so $\dimm\nul(B)\ge n-k$ (Rank--Nullity). Consider the
$(k+1)$-dimensional subspace $T=\spn\{\bv_1,\dots,\bv_{k+1}\}$. Since
\[
\dimm\nul(B)+\dimm T\ \ge\ (n-k)+(k+1)=n+1>n,
\]
the two subspaces of $\R^n$ intersect nontrivially; pick a unit vector
$\bw\in\nul(B)\cap T$. Write $\bw=\sum_{i=1}^{k+1}c_i\bv_i$ with $\sum_i c_i^2=1$. Then
$B\bw=\bzero$, so $(A-B)\bw=A\bw=\sum_{i=1}^{k+1}\sigma_ic_i\bu_i$, and by orthonormality
of the $\bu_i$,
\[
\norm{(A-B)\bw}^2=\sum_{i=1}^{k+1}\sigma_i^2c_i^2\ \ge\
\sigma_{k+1}^2\sum_{i=1}^{k+1}c_i^2=\sigma_{k+1}^2 ,
\]
using $\sigma_i\ge\sigma_{k+1}$ for $i\le k+1$. Therefore
$\sn{A-B}\ge\norm{(A-B)\bw}/\norm{\bw}\ge\sigma_{k+1}$.
\end{proof}

\begin{remark}[Weyl's inequality, optional]
The spectral result is a special case of \emph{Weyl's inequality}
$\sigma_{i+j-1}(X+Y)\le\sigma_i(X)+\sigma_j(Y)$: with $X=A-B$, $Y=B$, $j=k+1$, and
$\sigma_{k+1}(B)=0$ (since $\rank B\le k$), one gets
$\sigma_{i+k}(A)\le\sigma_i(A-B)$, whose $i=1$ case is the theorem and whose squares
summed give an alternative proof of the Frobenius version.
\end{remark}

\begin{ml_box}
The two norms answer two questions. Frobenius ($\sum_{i>k}\sigma_i^2$) is the
\emph{total} reconstruction energy lost --- the right metric for compression and PCA,
where average error matters. Spectral ($\sigma_{k+1}$) is the \emph{worst-case}
direction-wise error --- the right metric for stability and for guarantees of the form
``no single mode is mis-estimated by more than $\sigma_{k+1}$.'' Remarkably, the same
matrix $A_k$ is optimal for both.
\end{ml_box}

\section*{Summary}
\addcontentsline{toc}{section}{Summary}
The truncation $A_k$ has Frobenius error $\sqrt{\sum_{i>k}\sigma_i^2}$ and spectral error
$\sigma_{k+1}$ (Proposition~\ref{prop:truncerr}). Lemma~\ref{lem:projbest} converts
``best rank-$k$ matrix'' into ``best $k$-dimensional subspace''; Lemma~\ref{lem:kyfan}
(Ky Fan) solves that subspace problem; together they prove the Frobenius optimality of
$A_k$ (Theorem~\ref{thm:eyf}). A dimension-counting argument proves spectral optimality
(Theorem~\ref{thm:eys}). In both norms, $A_k$ is best, and the discarded singular values
quantify the error exactly.

\begin{exercise}
Prove that if $\sigma_k=\sigma_{k+1}$ the minimizer in Theorem~\ref{thm:eyf} need not be
unique, by exhibiting a matrix and two distinct optimal $A_k$'s.
\end{exercise}
\begin{exercise}
Deduce from Theorem~\ref{thm:eyf} that the relative error
$\fn{A-A_k}/\fn{A}=\sqrt{\sum_{i>k}\sigma_i^2/\sum_i\sigma_i^2}$ equals one minus the
captured ``energy fraction.''
\end{exercise}
\begin{exercise}
Using the intersection argument of Theorem~\ref{thm:eys}, prove the more general
statement $\sn{A-B}\ge\sigma_{k+1}$ holds even when $B$ ranges over all matrices with
$\rank(B)\le k$ that need not share $A$'s singular vectors.
\end{exercise}
\begin{exercise}
Prove the Frobenius theorem a second way from Weyl's inequality (state and assume Weyl).
\end{exercise}

%##################################################################
\chapter{Intuition After the Proof: Where the SVD Lives}\label{ch:intuition}
%##################################################################

\begin{intuition_box}
Now that optimality is \emph{proved}, we may indulge in intuition without risk of
circularity. Each application below is the Eckart--Young--Mirsky theorem wearing
different clothes: replace a matrix by its best low-rank approximation and read off the
consequences.
\end{intuition_box}

\section{Geometry: hyper-ellipsoids and dominant directions}
$A$ sends the unit sphere of $\R^n$ to an ellipsoid in $\R^m$ with semi-axis lengths
$\sigma_i$ along $\bu_i$ (Chapter~\ref{ch:svd}). Truncating to $A_k$ keeps the $k$
longest axes and flattens the rest --- the best $k$-dimensional ``shadow'' of the
ellipsoid, by Theorem~\ref{thm:eyf}.

\section{Principal Component Analysis}
Let $X\in\R^{m\times n}$ hold $m$ data points in rows; centre it to $X_c=X-\mathbf 1\bmu\T$.
The covariance $\Sigma_{\mathrm{cov}}=\tfrac1m X_c\T X_c$ is symmetric PSD, and by
Corollary~\ref{cor:evdlink} its eigenvectors are the right singular vectors $\bv_i$ of
$X_c$ with eigenvalues $\lambda_i=\sigma_i^2/m$. Thus:
\[
\textbf{principal directions}=\{\bv_i\},\qquad
\textbf{variance along }\bv_i=\sigma_i^2/m .
\]
Projecting onto $\spn\{\bv_1,\dots,\bv_k\}$ gives the $k$-dimensional representation of
least squared reconstruction error --- exactly Theorem~\ref{thm:eyf} applied to $X_c$.
PCA \emph{is} truncated SVD of the centred data.

\section{Image compression}
A grayscale image is a matrix of intensities. Its singular values typically decay
quickly, so $A_k$ with modest $k$ is visually indistinguishable from $A$ while storing
$k(m+n+1)$ numbers instead of $mn$. The error $\sqrt{\sum_{i>k}\sigma_i^2}$ is the
provably smallest achievable at rank $k$.

\section{Latent Semantic Analysis and word embeddings}
For a term--document matrix $A$ (entry $=$ frequency of term $i$ in document $j$), the
truncated SVD $A_k=U_k\Sigma_kV_k\T$ maps terms and documents into a shared
$k$-dimensional ``concept'' space. Synonymous terms, though never co-occurring, acquire
nearby rows in $U_k\Sigma_k$ because they load on the same singular directions. Modern
word-embedding methods are descendants of this idea: a low-rank factorization of a
(transformed) co-occurrence matrix.

\section{Recommendation systems and matrix completion}
A user--item matrix is modelled as approximately low rank: tastes are governed by a few
latent factors. Filling missing entries and truncating by SVD predicts unseen ratings;
the low-rank assumption is exactly the statement that $A\approx A_k$, whose optimality
Theorem~\ref{thm:eyf} guarantees among rank-$k$ models.

\section{LoRA: low-rank adaptation of large models}
Fine-tuning a pretrained weight matrix $W_0$ produces an update $\Delta W$. Empirically
$\Delta W$ is close to low rank, so LoRA stores it factored as $\Delta W=BA$ with
$B\in\R^{m\times k}$, $A\in\R^{k\times n}$, $k\ll m,n$. Training only $B,A$ replaces $mn$
parameters with $k(m+n)$ --- the same compression accounting as
Chapter~\ref{ch:ey}, deployed on gradient updates instead of images. The viability of
LoRA is the empirical claim that the best rank-$k$ approximation of the update loses
little, i.e.\ that the tail $\sum_{i>k}\sigma_i^2(\Delta W)$ is small.

\section{Why the SVD is central to machine learning}
Every task above reduces to the same operation: \emph{represent a matrix by its best
low-rank approximation}. The SVD provides that approximation, the theorem certifies its
optimality, and the singular values quantify exactly what is kept and what is lost. That
combination --- a universal decomposition plus a sharp optimality guarantee plus an
interpretable error --- is why the SVD recurs throughout statistics, signal processing,
and deep learning.

\section*{Summary}
\addcontentsline{toc}{section}{Summary}
PCA, compression, LSA, embeddings, recommenders, and LoRA are one theorem in disguise:
substitute $A_k$ for $A$. The right singular vectors are PCA's principal directions and
LSA's concepts; the singular values are variances, signal strengths, and compression
budgets; the discarded tail $\sum_{i>k}\sigma_i^2$ is the provably minimal information
lost.

\begin{exercise}
Show that the rank-$k$ PCA reconstruction $\hat X=\mathbf 1\bmu\T+(X_c V_k)V_k\T$ equals
$\mathbf 1\bmu\T+(X_c)_k$, the truncated SVD of the centred data, and conclude its
optimality from Theorem~\ref{thm:eyf}.
\end{exercise}
\begin{exercise}
For an image with singular values $\sigma_i\propto i^{-1}$, estimate the relative
Frobenius error of $A_k$ as a function of $k$.
\end{exercise}
\begin{exercise}
Argue that if a fine-tuning update has rapidly decaying singular values, a small LoRA
rank $k$ suffices, and quantify ``suffices'' via $\sum_{i>k}\sigma_i^2$.
\end{exercise}

\backmatter
\chapter{Notation and Theorem Dependency}
\addcontentsline{toc}{chapter}{Notation and Theorem Dependency}

\textbf{Notation.}
$\R^n$ real $n$-vectors; $A\T$ transpose; $\inner{\cdot}{\cdot}$ inner product;
$\norm{\cdot}$ Euclidean norm; $\fn{\cdot}$ Frobenius norm; $\sn{\cdot}$ spectral norm;
$\tr$ trace; $\col,\row,\nul$ the fundamental subspaces; $\rank$ rank; $P_S$ orthogonal
projector onto $S$; $\sigma_i$ singular values; $\bu_i,\bv_i$ left/right singular
vectors; $A_k$ rank-$k$ truncated SVD.

\medskip
\noindent\textbf{Dependency spine.}
Theorem~\ref{thm:rankequiv} (rank) and Theorem~\ref{thm:rank1} (rank-one)
$\to$ Theorem~\ref{thm:bestapprox} (best approximation) and
Corollary~\ref{cor:matpyth} (matrix Pythagoras)
$\to$ Theorem~\ref{thm:eigorth} (orthogonal eigenvectors)
$\to$ Theorem~\ref{thm:spectral} (Spectral Theorem)
$\to$ Propositions~\ref{prop:grampsd},~\ref{prop:sharedspec} (PSD)
$\to$ Theorem~\ref{thm:svdexist} (SVD exists)
$\to$ Theorem~\ref{thm:frobinv}, Corollary~\ref{cor:frobsing},
Theorem~\ref{thm:specnorm} (norms)
$\to$ Lemmas~\ref{lem:projbest},~\ref{lem:capture},~\ref{lem:kyfan}
$\to$ Theorems~\ref{thm:eyf},~\ref{thm:eys} (Eckart--Young--Mirsky).

\vspace{1cm}
\begin{center}
\color{steel}\itshape
``The SVD is the Spectral Theorem for rectangular matrices; Eckart--Young--Mirsky says
that keeping the largest singular values loses provably the least.''
\end{center}

\end{document}